\section{Anillos, homomorfismos, unidades y divisores de cero}\label{sec:anillos-fundamentos}

\subsection{Axiomas y primeros ejemplos}\label{tema:11-1-1}
\begin{defi}
Un \emph{anillo} es un conjunto $R$ con operaciones $+$ y $\cdot$ tales que $(R,+)$ es un grupo abeliano, el producto es asociativo y tiene identidad bilateral $1_R$, y se cumplen ambas leyes distributivas:
\[
r(s+t)=rs+rt,\qquad (r+s)t=rt+st.
\]
El neutro aditivo se denota $0_R$. El anillo es \emph{conmutativo} si $rs=sr$ para todos $r,s\in R$.
\end{defi}
Así, $(R,\cdot)$ es un monoide, pero no se pide que sus elementos tengan inversos multiplicativos.
En estas notas los anillos tienen unidad, pero no se supone conmutatividad salvo que se indique. Se permite el anillo cero, cuyo único elemento satisface $0_R=1_R$. Recíprocamente, si $0_R=1_R$, entonces $r=r1_R=r0_R=0_R$ para todo $r$, por la propiedad siguiente.

\begin{lema}\label{lema:C11-cero}
En cualquier anillo, $r0=0=0r$ y $(-1)r=-r=r(-1)$.
\end{lema}
\begin{proof}
La distributividad da $r0=r(0+0)=r0+r0$. Cancelando en el grupo aditivo resulta $r0=0$; el otro lado es análogo. Además,
\[
(-1)r+r=(-1+1)r=0,\qquad r(-1)+r=r(-1+1)=0.
\]
La unicidad del inverso aditivo concluye la prueba.
\end{proof}

\begin{ejem}
Los enteros, racionales, reales y complejos son anillos con sus operaciones usuales. También lo es $\mathbb Z/n\mathbb Z$ para $n\geq1$, con operaciones inducidas por las de los enteros. El caso $n=1$ es el anillo cero.

El anillo $M_2(\mathbb R)$ de matrices reales de tamaño dos no es conmutativo. Por ejemplo, tomemos
\[
A=\begin{pmatrix}0&1\\0&0\end{pmatrix},\qquad
B=\begin{pmatrix}0&0\\1&0\end{pmatrix}.
\]
Entonces $AB=\operatorname{diag}(1,0)$ y $BA=\operatorname{diag}(0,1)$ son distintos. La unidad es la matriz identidad.
\end{ejem}

\begin{ejem}[Fallo de cancelación]
En $\mathbb Z/6\mathbb Z$ se tiene $[2][4]=[2]=[2][1]$, aunque $[4]\ne[1]$. La diferencia se explica por $[2][3]=[0]$, con ambos factores distintos de cero.
\end{ejem}

\subsection{Divisores de cero y dominios}\label{tema:11-1-2}
\begin{defi}
Un elemento $a\in R$ es \emph{divisor de cero por la izquierda} si existe $b\ne0$ con $ab=0$, y \emph{por la derecha} si existe $b\ne0$ con $ba=0$.
\end{defi}
No se exige $a\ne0$. Así, en un anillo no nulo, el propio cero es divisor de cero por ambos lados. En un dominio se excluyen los divisores de cero no nulos.

\begin{prop}\label{prop:C11-cancelacion}
Un elemento $a$ no es divisor de cero por la izquierda si y sólo si $L_a:R\to R$, $x\mapsto ax$, es inyectiva. Para la derecha, el criterio es la inyectividad de $R_a:R\to R$, $x\mapsto xa$.
\end{prop}
\begin{proof}
Si $a$ no es divisor por la derecha y $xa=x'a$, entonces $(x-x')a=0$ fuerza $x=x'$. Recíprocamente, si $R_a$ es inyectiva y $xa=0=0a$, resulta $x=0$, de modo que no hay un testigo no nulo de que $a$ sea divisor por la derecha. Cambiando el lado de los productos se obtiene la primera equivalencia. Estas aplicaciones son homomorfismos de grupos aditivos; no se afirma que preserven la unidad ni que sean morfismos de anillos.
\end{proof}

\begin{defi}
Un \emph{dominio entero}, o simplemente \emph{dominio}, es un anillo conmutativo no nulo tal que
\[
ab=0\quad\Longrightarrow\quad a=0\ \text{o}\ b=0.
\]
\end{defi}
Por tanto, en un dominio se puede cancelar cualquier factor no nulo. Los anillos $\mathbb Z$, $\mathbb Q$, $\mathbb R$ y $\mathbb C$ son dominios; $\mathbb Z/6\mathbb Z$ no lo es.

\subsection{Unidades y campos}\label{tema:11-1-3}
\begin{defi}
$u$ es una \emph{unidad por la izquierda} si existe $v$ con $uv=1$, y una \emph{unidad por la derecha} si existe $w$ con $wu=1$. Es una \emph{unidad bilateral}, o \emph{unidad}, si satisface ambas condiciones.
\end{defi}
\begin{obs}
Con esta convención, una unidad por la izquierda tiene un inverso por la derecha, y viceversa. Las ecuaciones $uv=1$ y $wu=1$ especifican el lado del inverso.
\end{obs}

\begin{prop}\label{prop:C11-unidades}
Para $u\in R$ se cumplen las siguientes propiedades.
\begin{enumerate}[label=\textup{(\roman*)}]
\item Existe $v$ con $uv=1$ si y sólo si $L_u$ es sobreyectiva; existe $w$ con $wu=1$ si y sólo si $R_u$ lo es.
\item Si $uv=1$, entonces $R_u$ es inyectiva. Si $wu=1$, entonces $L_u$ es inyectiva.
\item Si hay inversos por ambos lados, coinciden y son únicos.
\item Las unidades forman un grupo multiplicativo, denotado $R^\times$ o $R^*$.
\end{enumerate}
\end{prop}
\begin{proof}
Si $uv=1$, cualquier $y\in R$ se escribe $y=u(vy)$, así que $L_u$ es sobreyectiva. Recíprocamente, la sobreyectividad permite elegir $v$ con $L_u(v)=1$. La prueba para $R_u$ es análoga.

Si $xu=x'u$ y $uv=1$, multiplicar por $v$ a la derecha da $x=x'$. Si $ux=ux'$ y $wu=1$, multiplicar por $w$ a la izquierda da lo mismo. Esto prueba (ii), y la \cref{prop:C11-cancelacion} traduce el resultado en ausencia de divisores de cero del lado correspondiente.

Si $wu=1=uv$, entonces $w=w(uv)=(wu)v=v$. Todo inverso bilateral debe coincidir con éste, que denotamos $u^{-1}$.

Finalmente, $1$ es unidad y, si $u,v$ son unidades, $uv$ tiene inverso $v^{-1}u^{-1}$. El inverso de una unidad también es unidad. La asociatividad se hereda de $R$, con lo que
\[
R^\times=\set{u\in R\given u\text{ es unidad bilateral}}
\]
es un grupo.
\end{proof}

\begin{defi}
Un \emph{anillo de división} es un anillo en el que todo elemento no nulo es unidad bilateral. Un \emph{campo} es un anillo conmutativo \emph{no nulo} en el que todo elemento no nulo es unidad.
\end{defi}
Bajo esta definición se permite el anillo cero como anillo de división; los campos, en cambio, se suponen siempre no nulos.

\begin{prop}\label{prop:C11-finito}
Todo campo es un dominio, y todo dominio finito es un campo.
\end{prop}
\begin{proof}
En un campo, si $a\ne0$ y $ab=0$, multiplicar por $a^{-1}$ da $b=0$. Para la segunda afirmación basta probar que un dominio finito es campo. Para $a\ne0$, las aplicaciones $L_a$ y $R_a$ son inyectivas por la \cref{prop:C11-cancelacion}; como el conjunto es finito, son sobreyectivas. La \cref{prop:C11-unidades} da inversos por ambos lados, luego $a$ es unidad.
\end{proof}
La finitud es esencial: $\mathbb Z$ es dominio, pero no campo, pues $2$ no tiene inverso entero.

\begin{ejem}[Unidades módulo un entero]
Para $n\geq1$,
\[
(\mathbb Z/n\mathbb Z)^\times
=\set{[m]_n\given \mcd(m,n)=1}.
\]
En efecto, la imagen de $L_{[m]}$ es el subgrupo aditivo generado por $[m]$. Es todo el grupo si y sólo si $\mcd(m,n)=1$, por el cálculo de subgrupos cíclicos. Equivalentemente, una identidad de Bézout $am+bn=1$ da el inverso $[a]$, y un inverso módulo $n$ produce esa identidad.

Para $n\geq2$, $\mathbb Z/n\mathbb Z$ es dominio si y sólo si es campo, y esto ocurre si y sólo si $n$ es primo. Si $n$ es primo, cada clase no nula es coprima con $n$ y es unidad. Si $n=ab$ con $1<a,b<n$, las clases no nulas $[a]$ y $[b]$ tienen producto cero. El caso $n=1$ se excluye por ser el anillo cero.
\end{ejem}

\subsection{Homomorfismos y la categoría de anillos}\label{tema:11-2-2a}
\begin{defi}
Un \emph{homomorfismo de anillos} $\varphi:R\to S$ preserva ambas operaciones y la unidad:
\[
\varphi(a+b)=\varphi(a)+\varphi(b),\qquad
\varphi(ab)=\varphi(a)\varphi(b),\qquad
\varphi(1_R)=1_S.
\]
Los anillos con estos morfismos forman la categoría $\Ring$.
\end{defi}
Las identidades y composiciones preservan las tres condiciones. Por ser homomorfismo aditivo, $\varphi$ también preserva cero y negativos.

\begin{prop}\label{prop:C11-inicial}
El anillo cero es terminal en $\Ring$, y $\mathbb Z$ es inicial. El anillo cero no es inicial.
\end{prop}
\begin{proof}
La única función $R\to0$ es un homomorfismo, incluida la condición sobre la unidad, porque $1_0=0$. Un morfismo $0\to R$ debería enviar el mismo elemento a $0_R$ y a $1_R$; sólo existe si $R=0$, de modo que $0$ no es inicial.

Para cualquier anillo $R$, definimos $\iota_R:\mathbb Z\to R$ por $\iota_R(n)=n1_R$, donde se usan sumas repetidas para $n>0$, cero para $n=0$ y negativos para $n<0$. La aditividad es la regla de los múltiplos enteros en un grupo abeliano. Por distributividad,
\[
(m1_R)(n1_R)=(mn)1_R.
\]
Para enteros positivos se obtiene sumando $m$ paquetes de $n$ unos; los casos negativos se siguen de la compatibilidad del producto con los inversos aditivos. También $\iota_R(1)=1_R$, por lo que es homomorfismo. Cualquier otro debe enviar $1$ a $1_R$ y, por aditividad, $n$ a $n1_R$.
\end{proof}

\begin{prop}
Los homomorfismos de anillos preservan unidades por cualquiera de los lados y, en particular, inducen homomorfismos $R^\times\to S^\times$.
\end{prop}
\begin{proof}
Si $uv=1_R$, entonces $\varphi(u)\varphi(v)=\varphi(1_R)=1_S$. Lo mismo ocurre con un inverso por el otro lado. Para unidades bilaterales, $\varphi(u^{-1})=\varphi(u)^{-1}$, y la compatibilidad con el producto ya es parte de los axiomas.
\end{proof}
\begin{ejem}
No se preserva necesariamente la ausencia de divisores de cero: $3\in\mathbb Z$ no es divisor de cero, pero su imagen bajo la proyección $\mathbb Z\to\mathbb Z/6\mathbb Z$ lo es, pues $[3][2]=[0]$ y $[2]\ne[0]$.
\end{ejem}
